\chapter{Conclusion}

This thesis set out to navigate autonomously on physical fields when neither satellite positioning nor a reliable embedded map can be taken for granted. The guiding idea has been to reconstruct the observation function of the field with Gaussian processes and to embed this reconstruction within a Bayesian, particle-based navigation filter, so that a single probabilistic language carries both the uncertainty of the state and the uncertainty of the map.

The foundations laid in the first part recalled the Bayesian filters that cope with the non-linear and multimodal densities of terrain-aided navigation -- from the Kalman filter to the regularised and Rao-Blackwellised particle filters -- and the Gaussian-process machinery -- kriging, its reduced-rank approximations and the reproducing-kernel-Hilbert-space viewpoint -- that turns field data into an observation function endowed with its own uncertainty. The second part brought the two together. Kriging was integrated into the particle filter, yielding a navigation scheme that no longer requires a dense and exact embedded map but reconstructs the field from whatever data are available, and that feeds the uncertainty of this reconstruction back into the state estimate. To keep the reconstruction affordable and to respect the non-stationarity of real fields, an adaptive kriging scheme restricts it to the data relevant to the current region, the resulting subsampling error being bounded theoretically. Finally, for the case of an altogether unknown field, localisation and mapping were carried out jointly -- through a Rao-Blackwellised filter and a reproducing-kernel-Hilbert-space reconstruction from a dictionary of basis functions -- and the learning bias inherent to feeding the map with the vehicle's own estimated positions was identified and characterised.

\textbf{TODO: bilan quantitatif -- résumer ici les résultats des campagnes (précision, robustesse, coût AK/statique), une fois les chiffres figés.}

\section{Limitations}

Several limitations temper these results. Adaptive kriging trades exactness for tractability: restricting the reconstruction to a local window rests on a local-stationarity assumption that has the status of a heuristic rather than of a guarantee, and the cost of the reconstruction, although reduced, remains higher than the mere reading of a stored map. In the map-free case, the learning bias was identified and characterised but not corrected; in the absence of a joint law over the state and the map, the joint estimate is moreover left without a reference against which its consistency could be certified. The observation function, finally, is treated as time-invariant, and the results rest on simulation rather than on data from a field campaign.

\section{Perspectives}

Several directions extend this work. On the side of the models, a spatially varying set of hyperparameters, fitted offline, would relax the local-stationarity heuristic on which adaptive kriging rests, and richer covariance families would widen the range of fields the method can handle. On the side of the methodology, the learning bias of the map-free case calls for a correction rather than a mere characterisation, and for a way of certifying the consistency of a joint estimate that no joint law underwrites.

The most natural opening, however, is the fusion of several complementary physical fields -- magnetic, gravimetric, bathymetric. Their joint informativeness could be quantified through the Fisher information, and their combination could resolve the multimodal likelihoods that a single field leaves ambiguous, extending the reconstruction pursued throughout this thesis from one field to several.
